% الجزء: sets-functions-relations
% الفصل: relations
% القسم: reflections
%
\documentclass[../../../include/open-logic-section]{subfiles}

\begin{document}

\olfileid[ar]{sfr}{rel}{ref}
\olsection{تأملات فلسفية}

في \olref[set]{sec}، عرّفنا العلاقات بأنها مجموعات معينة. وينبغي لنا أن
نتوقف قليلًا لنسأل سؤالًا فلسفيًا سريعًا: ما الذي \emph{يفعله} تعريف كهذا؟
من المشكوك فيه إلى حد بعيد أن نرغب في القول إننا \emph{اكتشفنا} بعض وقائع
الهوية الميتافيزيقية؛ كأن نقول، مثلًا، إن علاقة الترتيب على $\Nat$ قد
\emph{تبيّن أنها} المجموعة $R=  \Setabs{\tuple{n,m}}{n, m \in \Nat\text{ و }
n < m}$ التي عرّفناها في \olref[set]{sec}. وفيما يلي ثلاثة أسباب لذلك.

أولًا: في \olref[set][pai]{wienerkuratowski}، عرّفنا $\tuple{a, b} =
\{\{a\}, \{a, b\}\}$. وانظر بدلًا من ذلك في التعريف $\lVert a,
b\rVert = \{\{b\}, \{a, b\}\} = \tuple{b,a}$. عندما يكون $a \neq b$،
يكون لدينا $\tuple{a, b} \neq \lVert a,b\rVert$. لكن كان يمكننا بالقدر
نفسه أن نتخذ $\lVert a,b\rVert$ تعريفًا للزوج المرتب بدلًا من
$\tuple{a,b}$. وكان كلا التعريفين سيؤدي الغرض بالقدر نفسه. وهكذا أصبح
لدينا مرشحان صالحان بالقدر نفسه لأن يكون كل منهما «هو» علاقة الترتيب على
الأعداد الطبيعية، وهما:
\begin{align*}
		R &= \Setabs{\tuple{n,m}}{n, m \in \Nat \text{ و }n < m}\\
		S &= \Setabs{\lVert n,m\rVert}{n, m \in \Nat \text{ و }n < m}.
\end{align*}
وبما أن $R \neq S$، فإنه يتضح، بحكم الامتدادية، أنه لا يمكن أن يكونا
\emph{كلاهما} مطابقين لعلاقة الترتيب على~$\Nat$. لكن سيكون من التعسف
المحض، ومن ثم مما يدعو إلى شيء من الحرج، أن نزعم أن $R$ دون $S$ (أو
العكس) \emph{هي} علاقة الترتيب بحكم الواقع. (وهذه حالة بسيطة جدًا من حجة
ضد الاختزالية في نظرية المجموعات جعلها بيناسيراف مشهورة في
\citeyear{Benacerraf1965}. وسنعود إليها مرات عدة.)

ثانيًا: إذا رأينا أن \emph{كل} علاقة ينبغي أن تُماهى مع مجموعة، فإن علاقة
انتماء المجموعات نفسها، $\in$، ينبغي أن تكون مجموعة بعينها. بل يلزم أن
تكون المجموعة $\Setabs{\tuple{x,y}}{x \in y}$. ولكن هل هذه المجموعة
موجودة؟ نظرًا إلى مفارقة راسل، فإن وجود مثل هذه المجموعة ادعاء غير هيّن.
وفي الواقع، \oliflabeldef{cumul:::part}{فإن نظرية المجموعات التي نطوّرها في
\olref[cumul][][]{part} سوف \emph{تنكر} وجود هذه
المجموعة.\footnote{وباستباق ما سيأتي، إليك السبب. برهانًا بالخلف، افترض أن
$I = \Setabs{\tuple{x,y}}{x \in y}$ موجودة. عندئذ تكون $\bigcup \bigcup I$
هي المجموعة الشاملة، وهذا يناقض
\olref[sfr][z][sep]{thm:NoUniversalSet}.}}{يمكن تطوير نظرية المجموعات
تطويرًا صارمًا بوصفها نظرية بديهية، وستنكر تلك النظرية بالفعل وجود هذه
المجموعة.}
ومن ثم، حتى إذا أمكن معاملة بعض العلاقات كمجموعات، فلا بد أن تكون علاقة
انتماء المجموعات حالة خاصة.

ثالثًا: حين «ماهينا» العلاقات بالمجموعات، قلنا إننا سنجيز لأنفسنا كتابة
$Rxy$ بدلًا من $\tuple{x,y} \in R$. ولا بأس بهذا، شريطة أن تُعامل علاقة
الانتماء، «$\in$»، \emph{بوصفها} محمولًا. أما إذا اعتقدنا أن «$\in$» يرمز
إلى نوع معين من المجموعات، فإن العبارة «$\tuple{x,y} \in R$» لا تتألف إلا
من ثلاثة حدود مفردة تدل على مجموعات: «$\tuple{x,y}$» و«$\in$» و«$R$».
وليست قائمة كهذه من الأسماء أقدر على التعبير عن قضية من السلسلة عديمة
المعنى: «الكوب حامل الأقلام الطاولة». ومرة أخرى، حتى إذا
أمكن معاملة بعض العلاقات كمجموعات، فلا بد أن تكون علاقة انتماء المجموعات
حالة خاصة. (وهذا يجمع صيغة بسيطة من مفارقة مفهوم \emph{الحصان} عند فريغه
واعتراضًا شهيرًا كان فيتغنشتاين قد وجّهه إلى راسل.)

فإلام ينتهي بنا هذا؟ حسنًا، لا يوجد ما هو \emph{خطأ} في قولنا إن العلاقات
على الأعداد مجموعات. كل ما علينا هو أن نفهم الروح التي قيلت بها هذه
الملاحظة. فنحن لا نقرر واقعة هوية ميتافيزيقية. بل نلاحظ ببساطة أنه يمكننا،
في سياقات معينة، (وسنفعل) أن \emph{نعامل} علاقات (معينة) بوصفها مجموعات
معينة.

\end{document}
